\chapter{提问者人设库}
\label{app:personas}

提问者作为模拟用户，每个采样会话从人设库中\textbf{随机抽取一个人设并在该会话内保持固定}，以此替代真实用户在数据回流中的角色（机制见第~\ref{sec:cross_step} 节）。人设库共 42 个相互独立的人设，覆盖金融、审计、法务、运维、科研、教育、媒体、制造等职业场景。

每个人设由以下字段刻画：

\begin{itemize}
    \item \textbf{职业}与\textbf{偏好}：决定追问所关注的实质内容。如财务分析师要求每个数字可对账，运维工程师关心修复是否真正生效。
    \item \textbf{观察焦点}：两个二元维度的组合——\textbf{细节／整体}决定是追究单个数值还是审视交付物全貌，\textbf{内容／形式}决定是关注结果正确性还是关注排版与规范性。四种组合使不同人设从同一份观察报告中提取不同的追问依据。
    \item \textbf{语气}：平和／中性／急躁，仅影响表述风格，与耐心程度相互独立。
    \item \textbf{耐心} $P_0$ \textbf{与衰减系数} $d_0$：$P_0$ 为初始耐心值，每轮不满意后耐心乘以 $d_0$ 衰减，低于阈值即结束会话。放弃前的追问轮数约为 $\log_2(P_0 / d_0)$，故 $P_0$ 大、$d_0$ 接近 1 的人设会反复追问，$P_0$ 小、$d_0$ 低的人设则很快终止。该机制使会话长度随人设自然分化，而非由固定轮数上限统一截断。
\end{itemize}

全部 42 个人设及其参数如表~\ref{tab:personas} 所示。

\begin{longtable}{clp{3.1cm}cccc}
    \caption{提问者人设库（共 42 个人设）}
    \label{tab:personas} \\
    \hline
    编号 & 人设名称 & 职业 & 观察焦点 & 语气 & $P_0$ & $d_0$ \\
    \hline
    \endfirsthead
    \multicolumn{7}{l}{\small 续表~\thetable：提问者人设库} \\
    \hline
    编号 & 人设名称 & 职业 & 观察焦点 & 语气 & $P_0$ & $d_0$ \\
    \hline
    \endhead
    \hline
    \multicolumn{7}{r}{\small 续下页} \\
    \endfoot
    \hline
    \endlastfoot
        1 & Mei the financial analyst & Corporate financial analyst & 细节×内容 & 平和 & 5.7 & 0.63 \\
        2 & Greg the banker & Investment-bank executive & 整体×内容 & 急躁 & 7.1 & 0.71 \\
        3 & Omar the procurement specialist & Procurement specialist & 细节×内容 & 中性 & 5.9 & 0.64 \\
        4 & Hana the auditor & Auditor & 细节×内容 & 平和 & 9.3 & 0.86 \\
        5 & Diego the risk manager & Risk manager & 整体×内容 & 中性 & 5.2 & 0.59 \\
        6 & Raj the SRE & Site reliability / sysops engineer & 细节×内容 & 急躁 & 3.7 & 0.49 \\
        7 & Sven the sysadmin & System administrator & 细节×内容 & 中性 & 5.6 & 0.61 \\
        8 & Tariq the network engineer & Network engineer & 细节×内容 & 平和 & 7.4 & 0.74 \\
        9 & Bjorn the devops lead & DevOps lead & 整体×内容 & 中性 & 4.7 & 0.56 \\
        10 & Ken the compliance officer & Compliance / risk officer & 细节×内容 & 平和 & 3.3 & 0.47 \\
        11 & Carlos the project manager & Project manager & 整体×内容 & 中性 & 5.4 & 0.60 \\
        12 & Marcus the operations lead & Operations lead & 整体×内容 & 急躁 & 4.5 & 0.55 \\
        13 & Tom the startup founder & Early-stage startup founder & 整体×内容 & 急躁 & 3.2 & 0.46 \\
        14 & Lena the process consultant & Process-optimization consultant & 整体×内容 & 平和 & 9.5 & 0.87 \\
        15 & Felix the scheduler & Scheduling coordinator & 细节×内容 & 中性 & 7.3 & 0.72 \\
        16 & Nadia the project assistant & Project assistant & 整体×内容 & 平和 & 8.3 & 0.79 \\
        17 & Priya the support engineer & Customer support engineer & 细节×内容 & 中性 & 6.9 & 0.70 \\
        18 & Yusuf the call-center lead & Call-center team lead & 整体×内容 & 急躁 & 3.9 & 0.50 \\
        19 & Mira the conversation designer & Conversation experience designer & 整体×形式 & 平和 & 8.8 & 0.82 \\
        20 & Zane the sales director & Sales director & 整体×内容 & 急躁 & 3.5 & 0.48 \\
        21 & Sophie the content editor & Content / communications editor & 细节×形式 & 中性 & 6.6 & 0.68 \\
        22 & Yuki the UX writer & UX writer & 整体×形式 & 平和 & 8.6 & 0.81 \\
        23 & Elif the PR manager & PR manager & 整体×形式 & 急躁 & 3.0 & 0.45 \\
        24 & Pablo the copywriter & Marketing copywriter & 整体×形式 & 中性 & 5.0 & 0.58 \\
        25 & Grace the speechwriter & Speechwriter & 整体×形式 & 平和 & 9.7 & 0.88 \\
        26 & Aiko the trainer & Corporate trainer & 整体×形式 & 平和 & 8.5 & 0.80 \\
        27 & Bella the data analyst & Data analyst & 细节×内容 & 中性 & 8.1 & 0.78 \\
        28 & Dr. Lena the researcher & Academic researcher & 细节×内容 & 平和 & 9.1 & 0.85 \\
        29 & Arjun the intelligence analyst & Intelligence analyst & 细节×内容 & 平和 & 10.0 & 0.90 \\
        30 & Wei the research assistant & Research assistant & 细节×内容 & 中性 & 7.6 & 0.75 \\
        31 & Sara the consultant & Management consultant & 整体×内容 & 急躁 & 4.0 & 0.52 \\
        32 & Nina the QA tester & Quality assurance tester & 细节×内容 & 平和 & 9.8 & 0.89 \\
        33 & Igor the statistician & Statistician & 细节×内容 & 平和 & 9.0 & 0.83 \\
        34 & Tania the growth analyst & Growth analyst & 细节×内容 & 中性 & 6.4 & 0.67 \\
        35 & Hugo the experiment scientist & Experiment scientist & 细节×内容 & 中性 & 8.0 & 0.77 \\
        36 & Ravi the BI engineer & Business-intelligence engineer & 整体×内容 & 急躁 & 4.4 & 0.54 \\
        37 & Anna the office assistant & Executive office assistant & 整体×形式 & 平和 & 6.8 & 0.69 \\
        38 & Leo the generalist intern & Generalist intern & 整体×内容 & 中性 & 6.2 & 0.66 \\
        39 & Mona the admin lead & Administrative lead & 整体×形式 & 中性 & 4.9 & 0.57 \\
        40 & Theo the secretary & Secretary & 细节×形式 & 平和 & 7.8 & 0.76 \\
        41 & Carmen the HR specialist & HR specialist & 细节×形式 & 中性 & 6.1 & 0.65 \\
        42 & Boris the legal counsel & Legal counsel & 细节×内容 & 中性 & 4.2 & 0.53 \\
\end{longtable}

人设库的分布特征为：观察焦点以“细节×内容”（18 个）与“整体×内容”（13 个）为主，关注形式的人设（11 个）用于覆盖排版与规范类追问；语气分布为平和 16 个、中性 17 个、急躁 9 个；耐心 $P_0$ 的取值范围为 3.0$\sim$10.0，衰减系数 $d_0$ 的范围为 0.45$\sim$0.90。该分布使自产生任务在内容深度、形式规范与会话长度三个方向上均具备多样性，避免追问风格集中于单一模式。
